\documentclass{article}

\usepackage[english]{babel}
\usepackage[letterpaper,top=2cm,bottom=2cm,left=3cm,right=3cm,marginparwidth=1.75cm]{geometry}
\setlength{\parindent}{0pt}

\usepackage{amsmath}
\usepackage{amsthm}
\usepackage{amssymb}
\usepackage{graphicx}
\usepackage{float}
\usepackage{listings}
\usepackage[colorlinks=true, allcolors=blue]{hyperref}
\usepackage{xcolor}

\lstset{
  basicstyle=\ttfamily\small,
  numbers=left,
  numberstyle=\tiny,
  stepnumber=1,
  numbersep=8pt,
  backgroundcolor=\color{gray!10},
  frame=single,
  rulecolor=\color{gray!60},
  tabsize=2,
  captionpos=b,
  breaklines=true,
  breakatwhitespace=true,
  showstringspaces=false,
  keywordstyle=\color{blue},
  commentstyle=\color{green!50!black},
  stringstyle=\color{orange},
}

\title{DSC 291: Learning Theory Assignment 1}
\author{Tongle Shen}

\begin{document}
\maketitle

\section{Part B: Theory Problem}

\subsection{Problem 1}

\subsubsection{Choice of Grid}

Let
$$
m = \left\lceil \frac{1}{\Delta} \right\rceil
\quad \text{and} \quad
G = \left\{ \frac{k}{m} : k = 0,1,\dots,m \right\}.
$$
Then $G \subseteq [0,1]$ is a finite grid with
$$
|G| = m+1 = \left\lceil \frac{1}{\Delta} \right\rceil + 1.
$$

\subsubsection{Existence of a Consistent Grid Threshold}

We claim that for every $\Delta$-separated threshold-realizable sequence, there exists some
$\widetilde{\theta} \in G$ such that
$$
h_{\widetilde{\theta}}(x_t) = y_t
\quad \text{for all } t = 1,\dots,T.
$$

Since the spacing of the grid is $1/m \leq \Delta$, there exists a grid point
$\widetilde{\theta} \in G$ nearest to $\theta^*$ such that
$$
|\widetilde{\theta} - \theta^*| \leq \frac{1}{2m} \leq \frac{\Delta}{2}.
$$

Now fix any round $t$.

\textbf{Case 1: } If $y_t = +1$, then
$$
y_t = h_{\theta^*}(x_t) = \operatorname{sign}(x_t - \theta^*)
$$
and the $\Delta$-separation assumption implies
$$
x_t \geq \theta^* + \Delta.
$$
Hence
$$
x_t - \widetilde{\theta}
\geq (\theta^* + \Delta) - \widetilde{\theta}
\geq \Delta - |\widetilde{\theta} - \theta^*|
\geq \frac{\Delta}{2}
> 0.
$$
Therefore,
$$
h_{\widetilde{\theta}}(x_t) = \operatorname{sign}(x_t - \widetilde{\theta}) = +1 = y_t.
$$

\textbf{Case 2: } If $y_t = -1$, then similarly the realizability and separation assumptions imply
$$
x_t \leq \theta^* - \Delta.
$$
Thus
$$
x_t - \widetilde{\theta}
\leq (\theta^* - \Delta) - \widetilde{\theta}
\leq -\Delta + |\widetilde{\theta} - \theta^*|
\leq -\frac{\Delta}{2}
< 0,
$$
so
$$
h_{\widetilde{\theta}}(x_t) = \operatorname{sign}(x_t - \widetilde{\theta}) = -1 = y_t.
$$

Since this holds in both cases, we have shown that one hypothesis in
$$
H_G = \{ h_\theta : \theta \in G \}
$$
labels the entire sequence correctly.

\subsubsection{Mistake Bound Derivation}

Now run the Halving algorithm on the finite class $H_G$: on each round, predict by the majority
vote of the currently surviving hypotheses, and after a mistake eliminate all hypotheses that
predicted incorrectly. Because the sequence is realizable by $h_{\widetilde{\theta}} \in H_G$, the
Halving theorem implies that the total number of mistakes is at most
$$
\log_2 |H_G| = \log_2 |G|
= \log_2 \left( \left\lceil \frac{1}{\Delta} \right\rceil + 1 \right).
$$

\subsubsection{Conclusion}

So an explicit online learner is the Halving algorithm over the grid class $H_G$, with:

\begin{itemize}
\item choice of grid:
$$
G = \left\{ \frac{k}{m} : k = 0,1,\dots,m \right\},
\quad m = \left\lceil \frac{1}{\Delta} \right\rceil;
$$
\item size of the grid:
$$
|G| = \left\lceil \frac{1}{\Delta} \right\rceil + 1;
$$
\item resulting mistake bound:
$$
M \leq \log_2 \left( \left\lceil \frac{1}{\Delta} \right\rceil + 1 \right)
= O(\log(1/\Delta)).
$$
\end{itemize}

\subsection{Problem 2}

\subsubsection{Chosen Feature Map}

Take $d=2$ and define
$$
\phi(x) = (x,1) \in \mathbb{R}^2.
$$
For the true threshold $\theta^*$, define
$$
u^* = \frac{(1,-\theta^*)}{\sqrt{1+(\theta^*)^2}}.
$$
Then $u^*$ is a unit vector, and for every $x \in [0,1]$,
$$
\langle u^*, \phi(x) \rangle
= \frac{x-\theta^*}{\sqrt{1+(\theta^*)^2}}.
$$
Therefore,
$$
\operatorname{sign}(\langle u^*, \phi(x) \rangle)
= \operatorname{sign}(x-\theta^*)
= h_{\theta^*}(x).
$$
So this representation realizes the threshold function as a linear separator in $\mathbb{R}^2$.

\subsubsection{Margin Lower Bound}

Now consider any $\Delta$-separated threshold-realizable example $(x_t,y_t)$. Since
$$
y_t = \operatorname{sign}(x_t-\theta^*)
$$
and
$$
|x_t-\theta^*| \geq \Delta,
$$
we have
$$
y_t (x_t-\theta^*) = |x_t-\theta^*| \geq \Delta.
$$
Hence
$$
y_t \langle u^*, \phi(x_t) \rangle
= \frac{y_t (x_t-\theta^*)}{\sqrt{1+(\theta^*)^2}}
= \frac{|x_t-\theta^*|}{\sqrt{1+(\theta^*)^2}}
\geq \frac{\Delta}{\sqrt{1+(\theta^*)^2}}.
$$
Since $\theta^* \in [0,1]$, we have
$$
1+(\theta^*)^2 \leq 2,
$$
so
$$
y_t \langle u^*, \phi(x_t) \rangle \geq \frac{\Delta}{\sqrt{2}}.
$$
Thus the sequence is linearly separable with margin
$$
\gamma \geq \frac{\Delta}{\sqrt{2}}.
$$
In particular, we may take
$$
c = \frac{1}{\sqrt{2}}.
$$

\subsubsection{Norm Bound}

For every $x \in [0,1]$,
$$
\|\phi(x)\|
= \sqrt{x^2+1}
\leq \sqrt{2}.
$$
Therefore we may take
$$
R = \sqrt{2}.
$$

\subsubsection{Mistake Bound Derivation}

By the Perceptron theorem, the number of mistakes is at most
$$
\frac{R^2}{\gamma^2}.
$$
Using
$$
R = \sqrt{2}
\quad \text{and} \quad
\gamma \geq \frac{\Delta}{\sqrt{2}},
$$
we obtain
$$
M \leq \frac{(\sqrt{2})^2}{(\Delta/\sqrt{2})^2}
= \frac{2}{\Delta^2/2}
= \frac{4}{\Delta^2}.
$$

\subsubsection{Conclusion}

One valid choice is:
\begin{itemize}
\item feature map:
$$
\phi(x) = (x,1) \in \mathbb{R}^2;
$$
\item unit separator:
$$
u^* = \frac{(1,-\theta^*)}{\sqrt{1+(\theta^*)^2}};
$$
\item margin lower bound:
$$
\gamma \geq \frac{\Delta}{\sqrt{2}};
$$
\item norm bound:
$$
\|\phi(x)\| \leq \sqrt{2};
$$
\item resulting Perceptron mistake bound:
$$
M \leq \frac{4}{\Delta^2}
= O(1/\Delta^2).
$$
\end{itemize}

\subsection{Problem 3}

\subsubsection{Why This Does Not Contradict Part 1}

There is no contradiction between the continuous-threshold impossibility phenomenon from lecture
and the result of Part 1 because the two statements concern different families of sequences.

The impossibility phenomenon from lecture applies to arbitrary threshold-realizable sequences. In
that setting, examples may appear arbitrarily close to the true threshold $\theta^*$, so the learner
may need finer and finer resolution to determine which threshold is correct. In particular, the
continuous class $\{h_\theta : \theta \in [0,1]\}$ cannot be treated as a fixed finite class in the
worst case.

Part 1 adds the extra assumption of $\Delta$-separation. This changes the problem substantially:
every example satisfies
$$
|x_t-\theta^*| \geq \Delta.
$$
Because of this gap around $\theta^*$, two thresholds that are sufficiently close to each other are
indistinguishable on every admissible example. Equivalently, once we know the threshold up to
resolution on the order of $\Delta$, we know all labels in the sequence. That is why a grid of size
$O(1/\Delta)$ is enough in Part 1, and the Halving algorithm can be run on that finite effective
class.

So the lecture impossibility result rules out a uniform finite-mistake guarantee for the unrestricted
continuous threshold class, whereas Part 1 proves a finite-mistake guarantee on the smaller class of
$\Delta$-separated realizable sequences.

\subsubsection{Comparison of the Two Bounds}

The two bounds are
$$
M_{\text{Halving}} = O(\log(1/\Delta))
\quad \text{and} \quad
M_{\text{Perceptron}} = O(1/\Delta^2).
$$
They scale differently because they come from two different kinds of arguments.

\textbf{Part 1: } The Halving argument is combinatorial. After imposing $\Delta$-separation, the
continuous threshold class can be replaced by a finite grid of size $|G| = O(1/\Delta)$. The
Halving theorem then gives a mistake bound logarithmic in the number of candidate hypotheses:
$$
M \leq \log_2 |G| = O(\log(1/\Delta)).
$$
This argument measures how many distinguishable threshold locations remain once the margin gap is
taken into account.

\textbf{Part 2: } The Perceptron argument is geometric. It does not count hypotheses. Instead, it
uses the margin bound
$$
\gamma = \Omega(\Delta)
$$
and the norm bound
$$
R = O(1).
$$
The Perceptron theorem then yields
$$
M \leq \frac{R^2}{\gamma^2} = O(1/\Delta^2).
$$
This argument measures how far the examples are from the separating hyperplane in the chosen
linear representation.

\subsubsection{Interpretation}

The logarithmic bound from Part 1 is sharper for this particular problem because it exploits the
special one-dimensional threshold structure very directly. Once the threshold is localized to one of
only $O(1/\Delta)$ relevant positions, the problem is essentially finite.

The Perceptron bound is more general but also looser here. It applies through a generic margin
analysis for linear separators, and that analysis depends on $1/\gamma^2$. Since the margin is only
of order $\Delta$, the resulting guarantee is of order $1/\Delta^2$.

\subsubsection{Conclusion}

Part 1 and the lecture impossibility result are compatible because Part 1 assumes an additional
structural condition, namely $\Delta$-separation. Under that condition, the threshold problem has a
finite effective resolution and admits a logarithmic Halving bound. The Perceptron analysis gives a
different guarantee because it measures geometric margin rather than the number of distinguishable
thresholds, which is why it scales as $O(1/\Delta^2)$ instead of $O(\log(1/\Delta))$.

\end{document}
